% ============================================================
% Pictures of Inference - shared LaTeX preamble
% Imported by main.tex and any standalone chapter compile
% ============================================================

\documentclass[11pt,twoside,openany]{book}

% --- Geometry ---
\usepackage[
  letterpaper,
  inner=1.2in, outer=1.0in, top=1.0in, bottom=1.0in,
  marginparwidth=0pt, marginparsep=0pt
]{geometry}

% --- Fonts ---
\usepackage{mathpazo}
\usepackage[T1]{fontenc}
\usepackage[utf8]{inputenc}
\linespread{1.08}

% --- Math ---
\usepackage{mathtools}
\usepackage{amssymb}
\usepackage{bm}

% --- Headers ---
\usepackage{fancyhdr}
\pagestyle{fancy}
\fancyhf{}
\fancyhead[LE]{\small\itshape\nouppercase{\leftmark}}
\fancyhead[RO]{\small\itshape\nouppercase{\rightmark}}
\fancyfoot[C]{\thepage}
\renewcommand{\headrulewidth}{0.4pt}
\fancypagestyle{plain}{
  \fancyhf{}\fancyfoot[C]{\thepage}
  \renewcommand{\headrulewidth}{0pt}
}

% --- Chapter titling ---
\usepackage{titlesec}
\titleformat{\chapter}[display]
  {\normalfont\huge\bfseries}
  {\chaptertitlename\ \thechapter}{12pt}{\Huge}
\titlespacing*{\chapter}{0pt}{-20pt}{30pt}

% --- Graphics ---
\usepackage{graphicx}
% Graphics paths cover both build-from-prose and build-from-root scenarios
\graphicspath{
  {../pictures/figures/output/}%
  {../../pictures/figures/output/}%
  {pictures/figures/output/}%
  {figures/}%
}
\usepackage{wrapfig}
\usepackage{float}
\usepackage[font=small,labelfont=bf,labelsep=period]{caption}

% --- Code ---
\usepackage{listings}
\usepackage{xcolor}
\definecolor{codebg}{HTML}{f8f5f0}
\definecolor{codeframe}{HTML}{d4c9b0}
\definecolor{codekey}{HTML}{1a4f7a}
\definecolor{codecomment}{HTML}{7a7a7a}
\definecolor{codestr}{HTML}{b85c38}
\lstset{
  basicstyle=\small\ttfamily,
  keywordstyle=\bfseries\color{codekey},
  commentstyle=\itshape\color{codecomment},
  stringstyle=\color{codestr},
  backgroundcolor=\color{codebg},
  frame=single,
  rulecolor=\color{codeframe},
  breaklines=true,
  xleftmargin=8pt, xrightmargin=8pt,
  aboveskip=10pt, belowskip=10pt,
  columns=flexible,
  framesep=6pt,
}

\lstdefinelanguage{Rlang}{
  morekeywords={function,if,else,for,while,in,return,library,
    lm,summary,plot,hist,mean,sd,var,rnorm,set.seed,
    replicate,sample,data.frame,c,TRUE,FALSE,NULL,
    rbinom,ifelse,cat,str,head,nrow,ncol,sapply,class,
    cor,data,par,abline,curve,dnorm,t.test,sqrt,
    rexp,dpois,pnorm,qnorm,dbinom,polygon,seq,lines,glm,
    optim,solve,diag,binomial,IQR,median,summary,vif,anova,
    integrate,optimize,arima,arima.sim,acf,pacf,ivreg,subset,
    predict,exp,log,as.numeric,confint,coef,runif,points},
  morecomment=[l]{\#}, morestring=[b]", sensitive=true,
}

\lstdefinelanguage{Stata}{
  morekeywords={generate,gen,replace,drop,keep,merge,reshape,
    regress,reg,xtreg,ivregress,probit,logit,oprobit,
    if,by,sort,bysort,egen,collapse,summarize,sum,
    tabulate,tab,correlate,corr,predict,margins,
    program,syntax,version,clear,use,save,export,import,
    local,global,scalar,matrix,foreach,forvalues,while,
    return,exit,display,disp},
  morecomment=[l]{*}, morecomment=[l]{//},
  morestring=[b]", sensitive=true,
}

% --- Callout boxes ---
\usepackage[breakable,skins]{tcolorbox}

% Key Insight: the core idea (orange)
\newtcolorbox{keyinsight}{
  enhanced, breakable,
  colframe=orange!70!black, colback=orange!5,
  fonttitle=\sffamily\bfseries\small, title=Key Insight,
  boxrule=0.8pt, left=10pt, right=10pt, top=8pt, bottom=8pt,
  before skip=12pt, after skip=12pt,
}

% Try It Now: an exercise the reader does in real time (blue)
\newtcolorbox{tryitnow}{
  enhanced, breakable,
  colframe=blue!50!black, colback=blue!3,
  fonttitle=\sffamily\bfseries\small, title=Try It Now,
  boxrule=0.8pt, left=10pt, right=10pt, top=8pt, bottom=8pt,
  before skip=12pt, after skip=12pt,
}

% Inside the Math: deeper formal treatment (gold)
\newtcolorbox{inthemath}{
  enhanced, breakable,
  colframe=yellow!50!black, colback=yellow!4,
  fonttitle=\sffamily\bfseries\small, title=Inside the Math,
  boxrule=0.8pt, left=10pt, right=10pt, top=8pt, bottom=8pt,
  before skip=12pt, after skip=12pt,
}

% Code Lab: working code, usually R or Python (green)
\newtcolorbox{codelab}{
  enhanced, breakable,
  colframe=green!50!black, colback=green!3,
  fonttitle=\sffamily\bfseries\small, title=Code Lab,
  boxrule=0.8pt, left=10pt, right=10pt, top=8pt, bottom=8pt,
  before skip=12pt, after skip=12pt,
}

% Expedition: open-ended exploration with a real dataset (teal)
\newtcolorbox{expedition}{
  enhanced, breakable,
  colframe=teal!60!black, colback=teal!4,
  fonttitle=\sffamily\bfseries\small, title=Expedition,
  boxrule=0.8pt, left=10pt, right=10pt, top=8pt, bottom=8pt,
  before skip=12pt, after skip=12pt,
}

% Frontier Note: connection to current research (purple)
\newtcolorbox{frontiernote}{
  enhanced, breakable,
  colframe=purple!50!black, colback=purple!3,
  fonttitle=\sffamily\bfseries\small, title=Frontier Note,
  boxrule=0.8pt, left=10pt, right=10pt, top=8pt, bottom=8pt,
  before skip=12pt, after skip=12pt,
}

% --- Misc ---
\usepackage{booktabs}
\usepackage{enumitem}
\usepackage{microtype}
\usepackage{hyperref}
\hypersetup{
  colorlinks=true,
  linkcolor=blue!50!black,
  urlcolor=blue!50!black,
  citecolor=green!50!black
}

% --- Custom commands ---
\newcommand{\E}{\operatorname{\mathbb{E}}}
\newcommand{\Var}{\operatorname{Var}}
\newcommand{\Cov}{\operatorname{Cov}}
\newcommand{\Corr}{\operatorname{Corr}}
\newcommand{\bx}{\bm{x}}
\newcommand{\by}{\bm{y}}
\newcommand{\bbeta}{\bm{\beta}}
\newcommand{\bvarepsilon}{\bm{\varepsilon}}
\newcommand{\bX}{\bm{X}}
\newcommand{\bY}{\bm{Y}}
\newcommand{\bI}{\bm{I}}

% Convenient parsing aid: read-aloud notation
% Use \readas{notation}{spoken form} for first appearance
\newcommand{\readas}[2]{#1\;\textit{(read: #2)}}
